\documentclass[a4paper,12pt]{article}

% ══════════════════════════════════════════════
%   Pacotes
% ══════════════════════════════════════════════
\usepackage[utf8]{inputenc}
\usepackage[T1]{fontenc}
\usepackage[brazil]{babel}
\usepackage{helvet}
\renewcommand{\familydefault}{\sfdefault}

\usepackage{geometry}
\usepackage{microtype}
\usepackage{setspace}
\usepackage{enumitem}
\usepackage{booktabs}
\usepackage{array}
\usepackage{tabularx}
\usepackage{colortbl}
\usepackage{titlesec}
\usepackage{xcolor}
\usepackage{tcolorbox}
\usepackage{fancyhdr}
\usepackage{hyperref}
\usepackage{fontawesome5}
\usepackage{tikz}
\usepackage{calc}
\usepackage{graphicx}
\usepackage{ragged2e}

\tcbuselibrary{skins,breakable}

% ══════════════════════════════════════════════
%   Layout
% ══════════════════════════════════════════════
\geometry{
  left=2cm,
  right=2cm,
  top=2.3cm,
  bottom=2cm
}

% ══════════════════════════════════════════════
%   Cores
% ══════════════════════════════════════════════
\definecolor{corprimaria}{HTML}{0D2B45}
\definecolor{corsecundaria}{HTML}{1A5276}
\definecolor{coracento}{HTML}{2E86C1}
\definecolor{corcaixa}{HTML}{F0F6FC}
\definecolor{corlinha}{HTML}{BFD4E8}
\definecolor{corsuave}{HTML}{E8F0F8}
\definecolor{corverde}{HTML}{1D8348}
\definecolor{corlaranja}{HTML}{E67E22}
\definecolor{cortexto}{HTML}{2C3E50}
\definecolor{corcabec}{HTML}{0D2B45}

\color{cortexto}

\hypersetup{
  colorlinks=true,
  linkcolor=coracento,
  urlcolor=coracento
}

% Elimina warnings de \faIcon nos bookmarks do PDF
\pdfstringdefDisableCommands{%
  \def\faIcon#1{}%
  \def\hspace#1{ }%
  \def\raisebox#1#2{#2}%
}

% ══════════════════════════════════════════════
%   Tipografia e espaçamento
% ══════════════════════════════════════════════
\setstretch{1.15}
\setlength{\parindent}{0pt}
\setlength{\parskip}{4pt}
\setlength{\headheight}{14pt}

\setlist[itemize]{
  topsep=4pt,
  itemsep=2pt,
  leftmargin=1.3cm,
  label={\textcolor{coracento}{\faAngleRight}}
}
\setlist[enumerate]{
  topsep=4pt,
  itemsep=4pt,
  leftmargin=1.3cm
}

% ══════════════════════════════════════════════
%   Seções e subseções
% ══════════════════════════════════════════════
\titleformat{\section}
  {\Large\bfseries\color{corprimaria}}
  {}{0pt}{\raisebox{-1pt}{\rule{4pt}{14pt}\hspace{6pt}}}
  [\vspace{1pt}{\color{corlinha}\titlerule[0.8pt]}]
\titlespacing*{\section}{0pt}{18pt}{8pt}

\titleformat{\subsection}
  {\normalsize\bfseries\color{corsecundaria}}
  {}{0pt}{}
\titlespacing*{\subsection}{0pt}{10pt}{4pt}

% ══════════════════════════════════════════════
%   Caixas (tcolorbox)
% ══════════════════════════════════════════════
\newtcolorbox{blocosimples}{
  enhanced,
  colback=corcaixa,
  colframe=corlinha,
  arc=2mm,
  boxrule=0.4pt,
  left=10pt, right=10pt, top=10pt, bottom=10pt,
  shadow={0.5mm}{-0.5mm}{0mm}{black!8},
  breakable
}

\newtcolorbox{bloco}[2][]{
  enhanced,
  colback=white,
  colframe=corlinha!85!white,
  arc=2.2mm,
  boxrule=0.55pt,
  left=11pt, right=11pt, top=10pt, bottom=10pt,
  borderline west={1.6pt}{0pt}{coracento},
  before skip=10pt,
  after skip=10pt,
  fonttitle=\bfseries\color{white},
  title={\raisebox{-1pt}{#1}\hspace{4pt}#2},
  coltitle=white,
  attach boxed title to top left={xshift=0pt, yshift=-1.2mm},
  boxed title style={
    enhanced,
    colback=corprimaria!92!black,
    colframe=corprimaria!92!black,
    boxrule=0pt,
    arc=1.9mm,
    left=8pt, right=8pt, top=3.5pt, bottom=3.5pt
  },
  shadow={0.8mm}{-0.8mm}{0mm}{black!10},
  breakable
}

\newtcolorbox{blocobib}{
  enhanced,
  colback=white,
  colframe=corlinha,
  arc=2mm,
  boxrule=0.4pt,
  left=10pt, right=10pt, top=8pt, bottom=8pt,
  shadow={0.5mm}{-0.5mm}{0mm}{black!8},
  breakable
}

% ══════════════════════════════════════════════
%   Cabeçalho e Rodapé
% ══════════════════════════════════════════════
\pagestyle{fancy}
\fancyhf{}
\fancyhead[L]{\small\textcolor{corsecundaria}{Ementa 2026.1}}
\fancyhead[R]{\small\textcolor{corsecundaria}{Sistema de Cadastro em Python}}
\fancyfoot[C]{%
  \begin{tikzpicture}
    \fill[coracento] (0,0) circle (3.5mm);
    \node[text=white, font=\small\bfseries] at (0,0) {\thepage};
  \end{tikzpicture}%
}
\renewcommand{\headrulewidth}{0.35pt}
\renewcommand{\footrulewidth}{0pt}

% ══════════════════════════════════════════════
%   Comando para unidades do conteúdo
% ══════════════════════════════════════════════
\newcommand{\unidade}[3]{%
  \item \textbf{#1} \hfill
  \tikz[baseline=(n.base)]{\node[fill=coracento, text=white, font=\small\bfseries,
    rounded corners=2pt, inner sep=3pt, minimum width=2.2cm](n){#2};}%
  \begin{itemize}#3\end{itemize}%
}

% ══════════════════════════════════════════════
\begin{document}

% ──────────────────────────────────────────────
%   Título principal
% ──────────────────────────────────────────────
\begin{center}
  {\LARGE\bfseries\color{corprimaria}EMENTA DA UNIDADE CURRICULAR}\\[3pt]
  {\large\color{corsecundaria}Sistema de Cadastro em Python}
\end{center}

\vspace{0.2cm}

% ──────────────────────────────────────────────
%   Ficha de identificação
% ──────────────────────────────────────────────
\begin{tcolorbox}[
  enhanced,
  colback=white,
  colframe=corprimaria!80!black,
  arc=2.4mm,
  boxrule=0.55pt,
  left=8pt, right=8pt, top=7pt, bottom=7pt,
  shadow={0.7mm}{-0.7mm}{0mm}{black!10}
]
  {
    \setlength{\tabcolsep}{4pt}
    \renewcommand{\arraystretch}{1.34}
    \arrayrulecolor{corlinha!80!corprimaria}
    \begin{tabularx}{\linewidth}{
      @{}
      >{\RaggedRight\arraybackslash\bfseries\color{corprimaria}}m{4.6cm}
      >{\RaggedRight\arraybackslash}X
      >{\RaggedRight\arraybackslash\bfseries\color{corprimaria}}m{3.0cm}
      >{\RaggedRight\arraybackslash}m{2.8cm}
      @{}
    }
      \rowcolor{corsuave!85}
      \faIcon{book}\hspace{5pt}Unidade Curricular: &
      Sistema de Cadastro em Python &
      \faIcon{calendar-alt}\hspace{5pt}Semestre: &
      \textbf{2026.1} \\
      \specialrule{0.45pt}{0pt}{0pt}
      \faIcon{clock}\hspace{5pt}Carga Horária: &
      60 horas &
      \faIcon{cogs}\hspace{5pt}Natureza: &
      Teórico-prática \\
    \end{tabularx}
  }
\end{tcolorbox}

\vspace{0.3cm}

% ──────────────────────────────────────────────
%   Ementa
% ──────────────────────────────────────────────
\section{\faIcon{file-alt}\hspace{6pt}Ementa}

Definição e importância dos algoritmos, ferramentas e fluxogramas; problemas clássicos de lógica; definição e uso de variáveis, tipos de dados, operadores aritméticos, relacionais e lógicos; estruturas condicionais (\texttt{if}, \texttt{else}, \texttt{elif}) e de repetição (\texttt{for}, \texttt{while}); aninhamento de estruturas de controle; conceito e uso de funções, parâmetros, valores de retorno e modularização do código; manipulação de arquivos e persistência de dados; estruturas de dados (listas, tuplas, dicionários e conjuntos); desenvolvimento de um sistema de cadastro completo em Python; práticas de elaboração e utilização em programas.

% ──────────────────────────────────────────────
%   Objetivos
% ──────────────────────────────────────────────
\section{\faIcon{bullseye}\hspace{6pt}Objetivos}
Capacitar o(a) estudante a desenvolver, em Python, um sistema de cadastro simples e funcional, aplicando fundamentos de algoritmos e lógica de programação, estruturas de controle, funções e estruturas de dados, com foco em boas práticas de codificação e resolução de problemas.

%Capacitar o aluno a Compreender lógica de programação, aplicando conceitos de estruturas de controle, funções e manipulação de dados.


\begin{bloco}[\faIcon{list-ul}]{Objetivos Específicos}
\begin{itemize}
  \item Compreender os fundamentos de algoritmos e lógica de programação;
  \item Utilizar variáveis, tipos de dados e operadores em Python;
  \item Aplicar estruturas condicionais e de repetição na resolução de problemas;
  \item Desenvolver funções para modularização e reutilização de código;
  \item Manipular estruturas de dados como listas, tuplas e dicionários;
%   \item Implementar persistência de dados utilizando arquivos;
  \item Construir um sistema de cadastro funcional em Python.
\end{itemize}
\end{bloco}

% ──────────────────────────────────────────────
%   Conteúdo Programático
% ──────────────────────────────────────────────
\section{\faIcon{sitemap}\hspace{6pt}Conteúdo Programático}

\begin{blocosimples}
\begin{enumerate}[label=\textbf{Unidade \Roman*}, leftmargin=2.4cm, labelsep=0.4cm, itemsep=10pt]

  \unidade{Introdução à Lógica de Programação e Python}{8 horas}{
    \item Definição e importância dos algoritmos
    \item Ferramentas de desenvolvimento (google colab)
    \item Problemas clássicos de lógica
    \item Primeiro programa em Python
  }

  \unidade{Variáveis, Tipos de Dados e Operadores}{8 horas}{
    \item Definição e uso de variáveis
    \item Tipos de dados: \texttt{int}, \texttt{float}, \texttt{str}, \texttt{bool}
    \item Entrada e saída de dados (\texttt{input()}, \texttt{print()})
    \item Operadores aritméticos, relacionais e lógicos
    \item Conversão de tipos (casting)
  }

  \unidade{Estruturas Condicionais}{8 horas}{
    \item Estruturas \texttt{if}, \texttt{if/else} e \texttt{if/elif/else}
    \item Aninhamento de estruturas condicionais
    \item Exercícios práticos
  }

  \unidade{Estruturas de Repetição}{8 horas}{
    \item Estruturas \texttt{for} e \texttt{while}
    \item Comandos \texttt{break}, \texttt{continue} e \texttt{pass}
    \item Aninhamento de estruturas de repetição
    \item Exercícios práticos
  }

  \unidade{Funções e Modularização}{8 horas}{
    \item Conceito e definição de funções (\texttt{def})
    \item Parâmetros e argumentos
    \item Valores de retorno (\texttt{return})
    \item Escopo de variáveis (local e global)
    \item Funções com valores padrão
    \item Modularização do código
  }

  \unidade{Estruturas de Dados}{10 horas}{
    \item Listas: criação, acesso, métodos e iteração
    \item Tuplas: conceito e utilização
    \item Dicionários: chave-valor, métodos e iteração
    \item Conjuntos (\texttt{set}): operações básicas
    \item Lista de dicionários para cadastros
  }

  \unidade{Projeto: Sistema de Cadastro em Python}{Projeto Final}{
    \item Planejamento e modelagem do sistema
    \item Operações CRUD (Criar, Ler, Atualizar, Deletar)
    \item Validação de dados de entrada
    \item Apresentação do projeto final
  }

\end{enumerate}
\end{blocosimples}

\vspace{0.2cm}
\begin{center}
  \tikz{
    \node[fill=corprimaria, text=white, font=\bfseries, rounded corners=3pt,
      inner sep=6pt, minimum width=6cm]
      {\faIcon{clock}\hspace{6pt}Carga Horária Total: 60 horas};
  }
\end{center}

% ──────────────────────────────────────────────
%   Metodologia
% ──────────────────────────────────────────────
\section{\faIcon{chalkboard-teacher}\hspace{6pt}Metodologia}

\begin{bloco}[\faIcon{lightbulb}]{Estratégias de Ensino}
\begin{itemize}
  \item Aulas expositivas e dialogadas com auxílio de recursos audiovisuais;
  \item Resolução de exercícios práticos em laboratório de informática;
  \item Desenvolvimento de projetos individuais e em grupo;
  \item Atividades de fixação e avaliações práticas.
\end{itemize}
\end{bloco}

% ──────────────────────────────────────────────
%   Avaliação
% ──────────────────────────────────────────────
\section{\faIcon{chart-bar}\hspace{6pt}Avaliação}

\begin{bloco}[\faIcon{check-circle}]{Critérios e Instrumentos}
A avaliação será realizada de forma contínua, considerando:
\begin{itemize}
%   \item Participação e envolvimento nas atividades em sala;
  \item Exercícios e listas de atividades práticas;
  \item Avaliações individuais (provas);
  \item Projeto final: Sistema de Cadastro em Python.
\end{itemize}
\end{bloco}

% ──────────────────────────────────────────────
%   Bibliografia
% ──────────────────────────────────────────────
\section{\faIcon{book-open}\hspace{6pt}Bibliografia}

\subsection{\faIcon{bookmark}\hspace{4pt}Bibliografia Básica}
\begin{blocobib}
\begin{enumerate}[leftmargin=1.1cm, label={\textcolor{coracento}{\textbf{[\arabic*]}}}]
  \item FORBELLONE, A. L. V.; EBERSPACHER, H. F. \textbf{Lógica de programação}: a construção de algoritmos e estruturas de dados. São Paulo: Prentice Hall, 2005.
  \item GUEDES, S. \textbf{Lógica de programação algorítmica}. São Paulo: Pearson, 2014.
  \item XAVIER, G. F. C. \textbf{Lógica de programação}. São Paulo: Senac, 2014.
\end{enumerate}
\end{blocobib}

\subsection{\faIcon{bookmark}\hspace{4pt}Bibliografia Complementar}
\begin{blocobib}
\begin{enumerate}[leftmargin=1.1cm, label={\textcolor{coracento}{\textbf{[\arabic*]}}}]
  \item DEITEL, H. M. \textbf{C\#: como programar}. São Paulo: Pearson, 2003.
  \item DEITEL, H. M.; DEITEL, P. J. \textbf{C++: como programar}. São Paulo: Pearson, 2006.
  \item DEITEL, H. M.; DEITEL, P. J. \textbf{Java: como programar}. São Paulo: Pearson, 2010.
  \item FLANAGAN, D. \textbf{JavaScript: o guia definitivo}. Porto Alegre: Bookman, 2013.
  \item SOMMERVILLE, I. \textbf{Engenharia de software}. São Paulo: Pearson, 2011.
\end{enumerate}
\end{blocobib}

\end{document}
